\section{Introduction}
\label{sec:intro}

An assistant that only talks can be wrong in one way: it says
something false, and a person reads it. An assistant that \emph{acts}
is different. Consider the request that runs through this paper:

\begin{quote}\itshape
``Email the three customers whose invoices are overdue, and remind me
on Friday if any of them hasn't paid.''
\end{quote}

\noindent Carrying it out means reading a spreadsheet, opening a
mailbox, writing three messages, sending them under the person's name,
and acting again on Friday. Each step is an act in the world, and a
sent email cannot be recalled. The properties that make a language
model useful for this---it reads anything, writes fluently, and fills
gaps on its own---are the properties that make it dangerous to hold
the keys.

Five failures follow. A model \emph{fabricates}: ``I've sent the three
reminders'' is a sentence it can produce whether or not any mail left.
It is \emph{steered by what it reads}: text inside a document or an
email can pose as instructions, and no reliable defence against this
\emph{prompt injection} exists inside the model
itself~\cite{greshake2023indirect,perez2022ignore}. Given a person's
credentials, it \emph{oversteps}, because nothing stands between a
mistaken idea and its consequence. It \emph{leaves no trace} other than
its own account of what it did. And its tools are \emph{strangers'
code}, with their own dependencies, any of which may be careless or
hostile.

\paragraph{Approach}
\sys{} does not try to make the model trustworthy. It removes the
model from every decision that matters:

\begin{center}
\emph{The model may propose anything. It decides nothing.}
\end{center}

\noindent The model emits one structured action at a time---call this
function with these inputs, show this stored result, set this
reminder. Whether that action happens is decided by the platform: from
who the person is, what the chat permits, what the agent declared
before it was approved, and what the person approved, in code the
model cannot read, change, or argue with. None of these inputs is
derived from model output, so injected text can steer what the model
\emph{proposes} but cannot widen what may \emph{run}.

The design is deliberately unoriginal in its parts. It applies the
oldest ideas in operating systems~\cite{saltzer1975protection,
anderson1972planning} to a new kind of program: each agent is a
process running as its own user, its requests are system calls that
may be refused, its network access passes through a firewall, risky
actions ask the person first, a scheduler runs deferred work, a kill
switch ends what runs, and a log records what happened. What is new is
the component no operating system was designed for: a planner that is
highly capable and not trusted.

\paragraph{Contributions}
This paper makes the following contributions.

\begin{enumerate}[label=(\arabic*),leftmargin=1.6em]
\item A \textbf{separation of proposal from authority} for LLM agents,
  in which the model holds no permission, credential, or secret, and
  every proposed action is validated and gated by mechanisms none of
  whose inputs derive from model output (\cref{sec:design,sec:execution}).
\item A \textbf{deny-by-default authorization model} with two
  invariants checked on every change---an organization cannot lock
  itself out, and nobody can grant what they do not hold---and a
  \textbf{delegation} scheme under which the agent runtime acts for a
  person only through a short-lived, chat-bound, opaque token fenced to
  a fixed list of 40 actions (\cref{sec:authority}).
\item An \textbf{executor with call-bound approval}: an ordered
  sequence of gates on every function call, a person-chosen trust
  level, and approvals bound to the exact call by a hash that survives
  process restarts and is re-checked against current grants
  (\cref{sec:execution}).
\item \textbf{Evidence by construction}: anything displayed as data
  beside the model's words---tables, charts, files, ``verified''
  marks---is rendered from stored results of successful, schema-validated
  calls, never from the model's text (\cref{sec:evidence}).
\item A \textbf{manifest-first agent model}: each agent declares its
  functions, schemas, permission levels, resources, credential shapes,
  and network hosts before approval; approval pins one commit and one
  content-addressed package; and every call is held to the manifest
  (\cref{sec:agents}).
\item A \textbf{layered, self-inspecting sandbox} for untrusted agent
  code that needs no user namespaces: a capability-bearing spawn helper
  applies resource limits, \code{no\_new\_privs}, a per-agent uid, a
  Landlock ruleset and a seccomp denylist, behind an egress proxy and
  an nftables rule keyed on socket owner; the platform refuses to start
  where it cannot confine and reports what is not enforced where it
  runs (\cref{sec:sandbox}).
\item A \textbf{complete, open implementation} and an evaluation by
  more than two thousand automated tests, including a deliberately
  hostile agent, together with an explicit catalogue of limits
  (\cref{sec:implementation,sec:evaluation,sec:limitations}).
\end{enumerate}

\paragraph{What this paper does not claim}
\sys{} does not make a model correct, does not detect prompt injection,
and does not prevent an agent from misusing what it was rightly given.
Its claims are about containment and visibility: a mistake can lead
only to a call the person was already allowed to make, through an
agent the organization approved, reaching only declared places, with
anything above the chat's trust level waiting for the person's yes;
and every such call is recorded. \Cref{sec:limitations} lists where
these claims are weaker than they sound.

\paragraph{Organization}
\Cref{sec:threat} states the threat model. \Cref{sec:design} gives the
design and follows the running example end to end.
\Cref{sec:authority,sec:execution,sec:evidence,sec:agents,sec:sandbox,sec:secrets,sec:control}
describe each mechanism, naming the source files that implement it.
\Cref{sec:implementation} describes the implementation,
\cref{sec:evaluation} the evaluation, and \cref{sec:limitations} the
limits.
